\subsection{Roots in Lowest Terms} 
Sometimes, we want to rewrite a radical by putting a rational exponent in lowest terms. Usually, we will do this because we want to compare two results.
\begin{example}
Rewrite $\sqrt[6]{x^4}$ with the smallest possible index. Assume all variables represent positive real numbers.
\begin{lpic}{}{1}
\regtextleft{0.25}{-1}{$\sqrt[6]{x^4}=$};
\end{lpic}
\end{example}
By factoring them into primes, we can also do this with numbers.
\begin{example}
Rewrite $\sqrt[4]{25}$ with the smallest possible index.
\begin{lpic}{}{1}
\regtextleft{0.25}{-1}{$\sqrt[4]{25}=$};
\end{lpic}
\end{example}